\begin{abstract}
Tool-using LLM agents can commit security-relevant side effects before a human can inspect the realized operation. A safe pre-commit decision cannot depend only on the tool name, action text, or a coarse harmfulness label: it must bind the candidate action to the realized effect, target resource, operation mode, authorization context, and provenance/control source.

We define \emph{tool-effect binding} as this pre-commit safety property and build counterfactual stress tests that separate surface invariance from effect, resource, authorization, operation-mode, and provenance sensitivity. We then introduce effect-resource-operation atoms as the semantic object for mediation: a side-effectful tool call is a container of one or more atoms, each of which must be checked against explicit authorization and provenance context.

The resulting experiments show a consistent progression. On E48, a reference hard Effect-Binding Guard obtains 0.917 coverage with 0.036 unsafe pre-allow and 0.065 safe false-denial. On E50 resource/authorization stress, the same guard exposes the central bottleneck: 0.383 unsafe pre-allow at 0.921 coverage. E55-v2 then provides a controlled local warm-up: an authorization-aware atom-level prototype reaches 0.880 coverage with unsafe pre-allow 0/276 and safe false-denial 0/252.

We then test whether this evidence survives stronger objections. E60 introduces a 480-case independently specified held-out contract, where full atom-level mediation reaches 0.854 coverage with 0.000 unsafe pre-allow and 0.812 atom exact-set match. E61 evaluates 300 sandboxed realistic trace replays, reaching 0.847 coverage with 0.000 unsafe pre-allow, and adds a 156-trace saved AgentDojo-style/IPIGuard external replay subset with 0.904 coverage and 0.821 atom exact-set match. E62 decomposes extraction from authorization, E63 quantifies interface burden and context degradation, and E64 compares against stronger tool-call, resource-only, tuple-level, LLM-judge, least-privilege-style, and released-guardrail comparable baselines. These results establish counterfactual measurement and artifact-bounded evidence for atom-level pre-commit authorization, not a deployed-system guarantee.
\end{abstract}
